\begin{claim}{B.4}{Proved}{Accessibility for all computed coefficient theories}
For every $\U$-presentable enhanced coefficient category, and in
particular for the derivators of sets, groupoids, $r$-types ($r\geq0$), spaces,
$D(k)$ for a small associative unital ring $k$, spectra, and $L_E\Sp$
for a small spectrum $E$, the constant and $2$-complete classes are
set-generated fundamental localizers. They also satisfy the literal
accessibility condition Loc~(9) of Chapter XVI ms.~75--79: as full
subcategories of $\Cat^{[1]}$, they are generated under sufficiently
filtered colimits by a small family of bounded-cardinality arrows.
For the listed concrete coefficient theories the intermediate
classes are accessible full subcategories of
$\Cat^{[1]}$ as well; this does not turn them into fundamental
localizers when two-out-of-three fails. No large-cardinal principle
is used for these coefficient theories.
\end{claim}
\begin{proof}
We give the bounded-subfunctor construction, so that the two meanings
of accessibility are both verified. First fix $k\ne0$ and choose an
uncountable regular cardinal $\kappa\in\U$ larger than $|k|$ and
$\aleph_0$. Size of a category means its number of arrows, including
identities. Put $Q=\operatorname{Ex}^{\infty}$. The functor $Q$
preserves monomorphisms and filtered colimits. Every simplex of $QK$
occurs at a finite $\operatorname{Ex}$-stage and is a map from a
finite subdivided simplex; consequently it depends on finitely many
simplices of $K$. If a category has fewer than $\kappa$ arrows, its
$Q$-nerve has fewer than $\kappa$ simplices in total. For a vertex
$a$ of a Kan complex $K$, use the based-loop model
\[
 \Omega_aK=\{(a,a)\}\times_{K\times K}K^{\Delta^1}.
\]
The endpoint map is a Kan fibration. Each loop simplex uses only the
finite simplicial set $\Delta^n\times\Delta^1$, so this construction
also has finite support and commutes with filtered colimits of
pointed Kan complexes.

Let $f:A\to B$ belong to $W_{D(k)}^{\mathrm{lc}}$, and start with
any subfunctor $g_0:A_0\to B_0$ having fewer than $\kappa$ arrows.
Here $A_0,B_0$ are subcategories and $f(A_0)\subseteq B_0$.
We enlarge it inside $f$ using the criterion of C.7. For every
vertex $b\in B_0$, choose a vertex $a\in A$ and a finite zigzag
from $f(a)$ to $b$. For every pair $a,a'\in A_0$ whose images are
connected in $B_0$, choose a finite zigzag joining them in $A$.
Ambient bijectivity on components supplies these choices.

For every $a\in A_0$, every $n\geq0$, and every cycle
$z\in C_n(\Omega_{f(a)}QN B_0;k)$, choose an ambient cycle
$w\in C_n(\Omega_aQN A;k)$ and an ambient chain $b$ with
$f_*(w)-z=db$. Ambient surjectivity on loop homology supplies them.
For every cycle $w\in C_n(\Omega_aQN A_0;k)$ and every chain
$b\in C_{n+1}(\Omega_{f(a)}QN B_0;k)$ with $db=f_*(w)$, choose
an ambient chain $c$ with $dc=w$; ambient injectivity supplies it.
Each selected chain and zigzag has finite support. Pull the
supporting $Q$-nerve data back to finitely many original arrows of
$A$ and $B$, add those arrows and their images under $f$, and
generate subcategories by identities and finite composition.
There are fewer than $\kappa$ tests and witnesses: all relevant
chain groups have fewer than $\kappa$ elements because
$|k|<\kappa$. Thus the enlarged subfunctor $g_1$ still has size
less than $\kappa$.

Repeat this construction countably many times. Its union $g_\infty$
has size less than $\kappa$ by regular uncountability. Every
component relation in the union has a finite zigzag witness and
occurs at a finite stage. Every cycle and every bounding chain
likewise occurs at a finite stage, because its support is finite
and $QN$ and the loop model commute with filtered unions. At the
next stage the required lift or bounding chain was added. Thus
$g_\infty$ is a component bijection and its based-loop maps are
$k$-homology equivalences. C.7 gives
$g_\infty\in W_{D(k)}^{\mathrm{lc}}$.

Let $P_f$ be the poset of all size-$<\kappa$ subfunctors of $f$
lying in this class. Given fewer than $\kappa$ members, take their
union, generate subcategories, and apply the construction just
made. This gives an upper bound in $P_f$, so $P_f$ is
$\kappa$-filtered. Every original arrow lies in some member,
by starting the construction with that arrow. Therefore the
specified inclusions give
\[
 f=\colim_{g\in P_f}g\quad\text{in }\Cat^{[1]}.
\]
This also checks the enlargement form of the manuscript's Loc~(9),
with inclusions of source into source and target into target.
The same finite-support tests show closure of the class under
filtered colimits: $QN$, based loops with a chosen vertex, and
homology commute with them; bijectivity on components is detected
by finite zigzags. A category with fewer than $\kappa$ arrows is
$\kappa$-presentable in $\Cat$, and similarly for a functor in
$\Cat^{[1]}$: all its generators and composition equations form
a set of size $<\kappa$. These small weak arrows remain
$\kappa$-presentable in the full subcategory because its filtered
colimits are computed in the ambient arrow category. Choosing
representatives of their isomorphism classes therefore proves
accessibility as a full subcategory.

For set-generation as a fundamental localizer, let $S$ be that
small set of representatives and let $\langle S\rangle$ be its
generated fundamental localizer. B.2 gives
$\langle S\rangle\subseteq W_{D(k)}^{\mathrm{lc}}$.
Every $g\in P_f$ belongs to $\langle S\rangle$. Every
fundamental localizer is closed under filtered colimits, by
\citet[Corollary~4.2.22]{Cisinski2006}; hence
$f\in\langle S\rangle$. This proves equality. The cited
corollary alone is not an accessibility theorem; the bounded
construction is the indispensable extra argument.

For $L_E\Sp$, represent $E$ by a small cofibrant sequential
spectrum and choose $\kappa$ also larger than the total number
of simplices in all its levels. Use C.11 and replace loop-chain
cycles and boundaries by representatives and equality witnesses
in the stable homotopy groups of
$E\wedge\Sigma^\infty_+\Omega_aQN A$. A stable class is
represented at a finite spectrum level by a finite simplicial
sphere after levelwise fibrant replacement; an equality or
nullity is witnessed at a further finite level by a finite
homotopy. Thus each witness uses finitely many simplices of
the loop complex. For a size-$<\kappa$ complex there are fewer
than $\kappa$ such representatives and witnesses, over all
integer degrees and spectrum levels. The same countable
enlargement, filtered-union verification, and generated-localizer
argument therefore apply. This requires neither connectivity
of $E$ nor a smashing localization. Spectra are the case
$E=\mathbb S$.

For the constant classes, omit the based-loop construction and
use the ordinary chain or suspension-spectrum homology test
of C.7 or C.11. Their fundamental-localizer assertion is B.1.
For the intermediate classes use the homotopy-fibre test C.2
instead of based loops. A homotopy fibre of a Kan map is
constructed from a path object and a finite pullback, so it
has the same finite-support and filtered-colimit properties.
The resulting full subcategories are accessible, but the
set-generation conclusion as a fundamental localizer is not
asserted for them.

For sets, groupoids, and $r$-types, C.4 reduces the tests to
finitely many homotopy groups and a surjectivity condition
in the next degree. Representatives and homotopies are finite
sphere and disk data in Kan replacements, so the same
enlargement applies with any uncountable regular $\kappa$.
For spaces use every homotopy degree, still a countable family
of finite-support tests. These prove literal accessibility
and, for their constant and $2$-complete fundamental
localizers, set-generation. Zero coefficients give the whole
arrow category and are included by discarding all tests.
\paragraph{All presentable enhanced coefficients.}
We strengthen the concrete argument to every $\U$-presentable
$\infty$-category $\mathcal C$. Choose an uncountable regular
$\kappa\in\U$ for which $\mathcal C$ is
$\kappa$-compactly generated. The category $\PrL_\kappa$ of
$\kappa$-compactly generated presentable categories and
colimit-preserving functors preserving $\kappa$-compact objects
is itself $\U$-presentable:
\citet[Proposition~5.5.7.10]{EF_LurieHTT} identifies it with
$\Cat_\infty^{\mathrm{Rex}(\kappa)}$, and
\citet[Lemma~4.8.4.2]{EF_LurieHA} applies to the
$\U$-small family of representatives of the $\kappa$-small simplicial shapes.
By \citet[Proposition~5.5.7.6 and Notation~5.5.7.7]{EF_LurieHTT},
limits in its dual category $\operatorname{Pr}^{\mathrm R}_\kappa$ are
computed in categories. Consequently
\[
 F_{\mathcal C}:\mathcal S\longrightarrow\PrL_\kappa,
 \qquad X\longmapsto\Fun(X,\mathcal C),\qquad f\longmapsto f_!,
\]
is the copower functor $X\mapsto\colim_X\mathcal C$ and
preserves all small colimits. The natural direction here is
left Kan extension, not restriction; its adjoint is the
specified restriction $f^*$. In particular $F_{\mathcal C}$
is accessible. Restriction is an equivalence exactly when
its adjoint is, so $\W^{\mathrm{lc}}_{\mathcal C}$ is the
inverse image under $\operatorname{Arr}(F_{\mathcal C})$
of the full subcategory of equivalences. For an accessible
category $\mathcal E$, that full subcategory of
$\operatorname{Arr}(\mathcal E)$ is accessible: it is
equivalent to $\mathcal E$ by $e\mapsto\id_e$, and its
inclusion preserves filtered colimits.
\citet[Proposition~5.4.6.6]{EF_LurieHTT}, applied to the
resulting homotopy pullback, proves accessibility of this
arrow class of spaces.

To obtain the literal class of ordinary functors, apply
the same pullback theorem to the accessible functor
\[
 N(\operatorname{Arr}(\Cat_\U))\longrightarrow
       \operatorname{Arr}(\mathcal S),\qquad u\longmapsto Bu.
\]
It preserves filtered colimits because nerves and classifying
spaces do. Since the equivalence condition is full and replete,
this pullback is the nerve of the ordinary full subcategory
of arrows in $\W^{\mathrm{lc}}_{\mathcal C}$. Accessibility
of that nerve is ordinary accessibility: compactness is
detected by its discrete mapping spaces, and a filtered
infinity-category admits a cofinal ordinary filtered poset
\citep[Proposition~5.3.1.18]{EF_LurieHTT}.
Thus a $\U$-small set of functors generates the class under
sufficiently filtered colimits. It is closed under all
filtered colimits by B.2 and Cisinski's Corollary~4.2.22.
Its closure under all filtered colimits is therefore exactly
the class, proving literal Loc~(9). The generated-localizer
argument above proves set-generation as well.

For the constant class, choose a $\U$-small set $G$
generating $\mathcal C$ under colimits. The colimit-preserving
functor $X\mapsto(X\otimes g)_{g\in G}$ from spaces to
$\mathcal C^G$ inverts precisely $\W_{\mathcal C}$ by C.3:
the equivalence test is closed under colimits in the coefficient
object, so it suffices to test the generators. The identical
inverse-image and ordinary-nerve argument applies. This
strengthens all the concrete constant and locally constant
calculations without a large-cardinal hypothesis.
\end{proof}

The order printed on pairs of subcategories in Chapter XVI
ms.~78 reads $X_i\subset Y_i$ and $X_j\subset Y_j$.
It must read $X_i\subset X_j$ and $Y_i\subset Y_j$.
The former expression is not even typed when source and
target categories differ; the proof above uses the latter.

For an arbitrary derivator, accessibility is not proved here.
The bounded tests above depend on a fixed small ring, spectrum,
or truncation theory. Filtered-colimit closure by itself does
not imply set-generation. If one assumes Vop\v enka's
principle, every fundamental localizer is accessible
\citep[Scholie~4.2.17]{Cisinski2006}; that additional
set-theoretic hypothesis is not assumed in this paper.
